% ============================================================================
%  ANEXO A: Archivos de configuración del Servidor Web
% ============================================================================

% Este anexo se puede habilitar descomentando las líneas correspondientes
% en main.tex si se desea incluir material complementario.

\section*{Anexo A: Archivos de configuración}

A continuación se detallan las referencias a los archivos de configuración
utilizados para el despliegue del servidor web sobre Debian GNU/Linux 12:

\subsection*{A.1. Directivas globales de TLS y WebSocket}

Ubicadas en \texttt{/etc/nginx/nginx.conf} o en \texttt{/etc/nginx/conf.d/},
incluyen los parámetros de versiones de protocolo (TLS 1.2 y 1.3), gestión de
sesiones compartidas y el bloque \texttt{map} para la conmutación dinámica de
cabeceras WebSocket.

\subsection*{A.2. Configuración de Virtual Hosts y Proxies Inversos}

Definida en \texttt{/etc/nginx/sites-available/sudoers.conf} (activada mediante
enlace simbólico en \texttt{/etc/nginx/sites-enabled/}), que contiene todos los
bloques \texttt{server} para la redirección HTTP hacia HTTPS, el sitio principal
y los proxies hacia CUPS, Filebrowser, Portainer y Zabbix.

\subsection*{A.3. Certificados digitales SSL/TLS}

Generados mediante OpenSSL en \texttt{/etc/ssl/certs/web.pem} y
\texttt{/etc/ssl/private/web.key}, con permisos restrictivos y extensión
\textit{Subject Alternative Name} (SAN) para el dominio y sus subdominios.
